\documentclass{article}
\usepackage{amsthm}
\usepackage{amssymb}
\usepackage{url}
\usepackage{imprimatur}
% With hyperref, a block built with the public profile links to its Lean.
\usepackage{xcolor}
\definecolor{leansource}{RGB}{43,91,140}
\definecolor{crossref}{RGB}{55,65,85}
\usepackage[colorlinks,urlcolor=leansource,linkcolor=crossref,citecolor=crossref]{hyperref}
% The statements are claims: the proofs are not in this development.
\newtheorem{claim}{Claim}
% The note's environments, as imprimatur.toml's [aliases] declares them.
\leanblockalias{leanclaim}{claim}
\leanblockalias{leannotation}{}
\leanblockalias{leanprovenance}{}
\title{What the archive claims}
\author{}
\date{}
\begin{document}
\maketitle

On 20 November 2026 Levent Alp\"oge announced a proof of the Riemann
Hypothesis found with Claude, and released it as Lean alone: a four-page
note\footnote{\url{https://alpo.ge/remark.pdf}} and an archive of about
175\,000 Lean files, with no human-readable argument. The question everyone
asked first was not whether the proof is right but what, exactly, the Lean
proves. This note states the five theorems the archive claims, in the words
of their Lean statements. Every claim below is generated by imprimatur from a
Lean file that imports Mathlib and states the five theorems with
\texttt{sorry} for every proof: the file says what is claimed, and proves
nothing. The archive is not here, and nothing in this note says that its
proofs are correct.

The two halves of the question are checked by two tools. Lean's Comparator
(\texttt{lake comparator}) checks that the archive proves exactly the
statements of a challenge file like the one behind this
note.\footnote{\url{https://github.com/Deicyde/riemann-hypothesis} does
this for the five statements.} imprimatur checks that this note says exactly
what that file says: every claim is rendered from its Lean statement, and
every hand-written sentence tied to Lean with \verb|\leancheck| is checked
against the statement it describes. Neither tool reads the proof.

\section{The Riemann Hypothesis}

Mathlib defines the Riemann Hypothesis as a proposition, and the archive's
headline theorem proves that proposition. The claim is stated here with the
definition unfolded: a zero of $\zeta$ that is not one of the trivial zeros,
and is not the point $1$, has real part one half.

\begin{leanclaim}[claim, unfold]{Riemann.riemannHypothesis}{The Riemann Hypothesis}
\end{leanclaim}

Two clauses deserve a word. The trivial zeros are the negative even integers,
where $\zeta$ vanishes for reasons that have nothing to do with the
hypothesis. The clause $s \ne 1$ is there because Mathlib's
\texttt{riemannZeta} is a function on all of $\mathbb{C}$: at the pole $s = 1$
it takes an arbitrary value, so a statement about its zeros must set that
point aside \leancheck[claim]{Riemann.riemannHypothesis}. The same claim,
written with its hypotheses as an implication rather than inside Mathlib's
definition, reads the same:

\begin{leanclaim}[claim]{Riemann.riemannZeta_eq_zero_re}{The Riemann Hypothesis, restated}
\end{leanclaim}

\section{Dirichlet $L$-functions}

The archive extends the method to the Generalized Riemann Hypothesis. A
Dirichlet character modulo $N$ is a homomorphism from the units of
$\mathbb{Z}/N$ to $\mathbb{C}$, extended by zero, and its $L$-function
$L(\chi, s)$ generalises $\zeta$, which is $L(1, s)$ for the trivial
character modulo~$1$. The first claim is that an $L$-function has no zero to
the right of the critical line. Its one exception is the trivial character at
$s = 1$, where $L(1, s)$ is $\zeta(s)$ and has its pole, which Mathlib again
gives an arbitrary value \leancheck[claim]{Riemann.LFunction_ne_zero_of_half_lt_re}.

\begin{leanclaim}[claim]{Riemann.LFunction_ne_zero_of_half_lt_re}{No zeros right of the line}
\end{leanclaim}

The Generalized Riemann Hypothesis itself is the next claim: every zero of an
$L$-function in the critical strip lies on the critical line. Mathlib's
characters are indexed by a modulus $N \in \mathbb{N}$ that must not be zero,
a hypothesis Lean keeps in an instance binder and the claim states in words.

\begin{leanclaim}[claim]{Riemann.LFunction_eq_zero_re}{The Generalized Riemann Hypothesis}
\end{leanclaim}

A character is primitive when it is not induced from a character of a smaller
modulus. The last claim is the hypothesis for primitive characters, which
implies the general one, since the $L$-function of an imprimitive character
differs from that of the primitive character inducing it by finitely many
Euler factors, none of which vanishes in the strip
\leancheck[claim]{Riemann.LFunction_eq_zero_re_of_isPrimitive, Riemann.LFunction_eq_zero_re}.

\begin{leanclaim}[claim]{Riemann.LFunction_eq_zero_re_of_isPrimitive}{The hypothesis for primitive characters}
\end{leanclaim}

\section{What is checked, and what is not}

The table below lists each claim with its Lean declaration and the axioms it
rests on. Every row says \textbf{sorry}: the proofs are not in this
development. A reader who wants more than the statements has two things to
do, neither of which this note can do for them: run the Comparator on the
archive against these statements, and read the archive.

\begin{leanprovenance}{Provenance}{Provenance}
\end{leanprovenance}

The notation the claims are written in is listed below, each entry with the
Lean it stands for. The entries are where the words of this note meet
Mathlib's definitions; in particular, $\zeta$ is Mathlib's
\texttt{riemannZeta}, whose value at $1$ is the reason for the clause
$s \ne 1$ above.

\begin{leannotation}{Notation}{Index of notation}
\end{leannotation}

\end{document}
